\chapter{10\,000 以内的数}\label{ch:g3:numbers}

一千、二千、三千\,\dots\ 这些大数，
都是利用一个巧妙的想法构造出来的：数字所在的\emph{数位}
决定了它表示的值。本章学习如何读数、写数、
比较和排列 $10\,000$ 以内的数。

\section{读数和写数}

\begin{definition}[数位值]\label{def:g3:numbers:place}
一个数用数字 $0, 1, 2, \dots, 9$ 书写。从右往左，
数位依次是：个位、十位、百位、千位。在
$3\,254$ 中：
\[
3\,254 = 3 \text{ 个千} + 2 \text{ 个百} + 5 \text{ 个十}
+ 4 \text{ 个一}.
\]
\end{definition}

\begin{omfigure}
\begin{tikzpicture}[scale=0.9]
  \draw (0,0) grid[xstep=2.3, ystep=0.85] (9.2,1.7);
  \node[font=\small] at (1.15,1.275) {千位};
  \node[font=\small] at (3.45,1.275) {百位};
  \node[font=\small] at (5.75,1.275) {十位};
  \node[font=\small] at (8.05,1.275) {个位};
  \node[font=\large, omDef] at (1.15,0.425) {$3$};
  \node[font=\large, omDef] at (3.45,0.425) {$2$};
  \node[font=\large, omDef] at (5.75,0.425) {$5$};
  \node[font=\large, omDef] at (8.05,0.425) {$4$};
\end{tikzpicture}

{\small $3\,254$ 的数位表，读作“三千二
百五十四”。}
\end{omfigure}

\begin{example}[占位的零]\label{ex:g3:numbers:zero}
“四千零七”写作 $4\,007$：千位上是 $4$，
个位上是 $7$，中间的\emph{零}用来占位，
使每个数字都在正确的数位上。没有它们，$47$ 就是
另一个完全不同的数！
\end{example}

\begin{example}[相邻数位的换算]\label{ex:g3:numbers:exchange}
10 个一是 1 个十；10 个十是 1 个百；10 个百是
1 个千。所以 $10$ 个百 $= 1\,000$，$32$ 个十
$= 320$：数 $320$ 里含有 $32$ 个十（而不只是数字
$2$！）。
\end{example}

\section{比较与排序}

\begin{method}[比较两个数]\label{met:g3:numbers:compare}
\begin{enumerate}
  \item 数一数位数：位数多的数更大
        （$1\,002 > 998$）。
  \item 位数相同时，逐位比较，
        从\emph{左}边数起，第一个不同的数字决定大小：
        $5\,614 > 5\,589$，因为百位上 $6 > 5$。
  \item 用符号 $<$（小于）或 $>$
        （大于）表示结果：开口总是朝着较大的数。
\end{enumerate}
\end{method}

\begin{omfigure}
\begin{tikzpicture}[scale=1.1]
  \draw[-Stealth] (-0.3,0) -- (10.5,0);
  \foreach \i in {0,1,...,10} \draw (\i,-0.08) -- (\i,0.08);
  \node[below, font=\small] at (0,-0.1) {$0$};
  \node[below, font=\small] at (5,-0.1) {$500$};
  \node[below, font=\small] at (10,-0.1) {$1000$};
  \fill[omThm] (2.4,0) circle (2pt) node[above, font=\small] {$240$};
  \fill[omDef] (6.9,0) circle (2pt) node[above, font=\small] {$690$};
  \fill[omMeth] (9.05,0) circle (2pt) node[above, font=\small] {$905$};
\end{tikzpicture}

{\small 数轴按百划分刻度，越靠右，
数越大：$240 < 690 < 905$。}
\end{omfigure}

\begin{example}\label{ex:g3:numbers:order}
按从小到大的顺序排列：$780$、$8\,700$、$807$、$87$。
先看位数：$87$（两位），再是 $780$ 和 $807$
（三位），最后是 $8\,700$（四位）。比较 $780$ 和 $807$，
看百位数字，$7 < 8$，所以 $780 < 807$。排好的顺序为：
\[
87 < 780 < 807 < 8\,700 .
\]
\end{example}

\section{奇数和偶数}

\begin{definition}[偶数、奇数]\label{def:g3:numbers:evenodd}
一个数若能分成两个相等的整数部分，则称它为\emph{偶数}\index{偶数}
——它的个位数字是 $0$、$2$、$4$、$6$ 或
$8$。否则称为\emph{奇数}\index{奇数}（个位数字为 $1$、
$3$、$5$、$7$ 或 $9$）。
\end{definition}

\begin{example}
$46$ 是\omterm{def:g3:numbers:evenodd}{偶数}（$46 = 23 + 23$）；$47$ 是\omterm{def:g3:numbers:evenodd}{奇数}（不能分成两份相等的整数，
否则需要半个物体）。只需看\emph{最后一位}数字：$3\,578$ 是
\omterm{def:g3:numbers:evenodd}{偶数}，$2\,341$ 是\omterm{def:g3:numbers:evenodd}{奇数}。
\end{example}

\begin{omfigure}
\begin{tikzpicture}[scale=0.45]
  \foreach \i in {0,...,5} \fill[omDef] ({\i},0) circle (0.28);
  \foreach \i in {0,...,5} \fill[omDef] ({\i},1) circle (0.28);
  \node[below, font=\small] at (2.5,-0.7) {$12$：偶数，两行数量相等};
  \begin{scope}[xshift=10.5cm]
  \foreach \i in {0,...,4} \fill[omThm] ({\i},0) circle (0.28);
  \foreach \i in {0,...,4} \fill[omThm] ({\i},1) circle (0.28);
  \fill[omThm] (5,0) circle (0.28);
  \draw[omExo, thick] (5,0) circle (0.5);
  \node[below, font=\small] at (2.5,-0.7) {$11$：奇数，多出一个单独的点};
  \end{scope}
\end{tikzpicture}

{\small \omterm{def:g3:numbers:evenodd}{偶数}可以平均分成两行；\omterm{def:g3:numbers:evenodd}{奇数}总会
多出一个单独的点。}
\end{omfigure}

\section{练习}

\begin{exercise}[$\star$]\label{exo:g3:numbers:1}
用数字写出：“两千三百四十五”；“六千零五十”；
“九千零一”。
\end{exercise}

\begin{exercise}[$\star$]\label{exo:g3:numbers:2}
用文字写出：$1\,208$；\quad $4\,070$；\quad $9\,999$。
\end{exercise}

\begin{exercise}[$\star$]\label{exo:g3:numbers:3}
在 $7\,382$ 中，十位上的数字是几？千位上的
数字是几？这个数一共含有多少个\emph{百}
（注意：不只是百位上的数字，可参考
\cref{ex:g3:numbers:exchange}）？
\end{exercise}

\begin{exercise}[$\star$]\label{exo:g3:numbers:4}
仿照\cref{def:g3:numbers:place}，将下列各数分解：
$5\,631$；\quad $2\,046$；\quad $8\,500$。
\end{exercise}

\begin{exercise}[$\star$]\label{exo:g3:numbers:5}
抄写下列各式，并在空处填入 $<$ 或 $>$：
\[
456 \;?\; 465, \qquad
1\,203 \;?\; 989, \qquad
6\,090 \;?\; 6\,009, \qquad
9\,999 \;?\; 10\,000 .
\]
\end{exercise}

\begin{exercise}[$\star$]\label{exo:g3:numbers:6}
从小到大排列：$530$；\ $3\,500$；\ $353$；\ $5\,033$；
\ $503$。
\end{exercise}

\begin{exercise}[$\star$]\label{exo:g3:numbers:7}
在一条以百为单位刻度、从 $0$ 到 $1\,000$ 的数轴上，
大致标出数 $150$、$480$、$820$ 和 $990$ 的位置。
\end{exercise}

\begin{exercise}[$\star$]\label{exo:g3:numbers:8}
是\omterm{def:g3:numbers:evenodd}{偶数}还是\omterm{def:g3:numbers:evenodd}{奇数}？\ $34$；\ $57$；\ $70$；\ $2\,481$；\ $6\,548$。
\end{exercise}

\begin{exercise}[$\star$]\label{exo:g3:numbers:9}
按步长数一数：从 $2\,970$ 出发，每次加 $10$，数五步；从 $860$ 出发，
每次加 $100$，数四步；从 $9\,996$ 出发，每次加 $1$，数四步。
数到最后一个时，你发现了什么？
\end{exercise}

\begin{exercise}[$\star\star$]\label{exo:g3:numbers:10}
用数字 $3$、$7$、$0$、$5$ 各恰好用一次，写出
可能的最大四位数，再写出最小的四位数（一个数
不能以 $0$ 开头）。
\end{exercise}

\begin{exercise}[$\star\star$]\label{exo:g3:numbers:11}
我是一个三位数。我的数字是 $2$、$5$ 和 $8$，我是\omterm{def:g3:numbers:evenodd}{偶数}，
并且比 $800$ 大。我是谁？
\end{exercise}
